\chapter{Couche 5 — L'accès aux machines louées}

C'est la surface la plus exposée de la plateforme : des ports ouverts, des
comptes \cmd{root}, et des utilisateurs qu'on ne connaît pas.

% -----------------------------------------------------------------------------
\mesure{Le mot de passe à usage unique}{mesure 14 · \fichier{services.py}}

\menace
Un mot de passe \cmd{root} par défaut, ou dérivé du nom de la machine, serait
le même partout : il suffirait d'en louer une pour entrer dans toutes les
autres.

\parade
Seize caractères tirés au sort par machine, affichés \textbf{une seule fois},
puis conservés chiffrés. L'utilisateur peut le régénérer à chaud sans détruire
sa machine.

\danslecode
Le mot de passe est injecté au lancement (\cmd{-e ROOT\_PASSWORD=...}), jamais
inscrit dans l'image. Le \fichier{Dockerfile} le dit explicitement :

\begin{lstlisting}[language=bash]
# ROOT_PASSWORD n'est VOLONTAIREMENT pas defini ici : une valeur inscrite dans
# l'image serait lisible par `docker inspect` et identique sur toutes les
# machines. La plateforme l'injecte au lancement.
\end{lstlisting}

Et si la plateforme oublie de le fournir, l'\fichier{entrypoint.sh} n'applique
pas de valeur de repli connue --- il en tire une au hasard :

\begin{lstlisting}[language=bash]
if [ -z "${ROOT_PASSWORD:-}" ]; then
    ROOT_PASSWORD="$(tr -dc 'A-Za-z2-9' < /dev/urandom | head -c 16)"
fi
\end{lstlisting}

\begin{retenir}
Une valeur par défaut est une valeur publique. Entre « échouer » et « retomber
sur une valeur connue », un système sûr échoue --- ou, comme ici, tire au
hasard et le journalise.
\end{retenir}

\limite
Le mot de passe transite par une page web et le serveur doit le connaître pour
l'afficher. C'est inhérent, et c'est la raison d'être de la mesure suivante.

% -----------------------------------------------------------------------------
\mesure{L'authentification par clé publique}{mesure 15 · \fichier{entrypoint.sh}}

\menace
Un port SSH ouvert sur Internet reçoit des milliers de tentatives par jour. Un
mot de passe de seize caractères y résiste --- mais la tentative reste
\emph{possible}, et le secret doit voyager pour être utilisé.

\parade
L'utilisateur dépose sa clé publique sur son compte. Les machines qu'il loue
\textbf{ensuite} la reçoivent, et leur \cmd{sshd} est reconfiguré pour refuser
purement et simplement la méthode « mot de passe ».

\danslecode
\begin{lstlisting}[language=bash]
if [ -n "$SSH_PUBKEY" ]; then
    mkdir -p /root/.ssh && chmod 700 /root/.ssh
    printf '%s\n' "$SSH_PUBKEY" > /root/.ssh/authorized_keys
    chmod 600 /root/.ssh/authorized_keys
    sed -i 's/^#\?PasswordAuthentication .*/PasswordAuthentication no/' \
        /etc/ssh/sshd_config
    sed -i 's/^#\?PermitRootLogin .*/PermitRootLogin prohibit-password/' \
        /etc/ssh/sshd_config
fi
\end{lstlisting}

Vérifié sur une machine louée, en forçant la méthode mot de passe :

\begin{lstlisting}[language=bash]
$ ssh -o PubkeyAuthentication=no -o PreferredAuthentications=password \
      root@127.0.0.1 -p 8096
root@127.0.0.1: Permission denied (publickey).
\end{lstlisting}

\begin{notion}{Le vrai gain : il n'y a plus de secret partagé}
Avec un mot de passe, le secret existe \textbf{à deux endroits} : l'utilisateur
le connaît, et le serveur doit le connaître pour le vérifier. Quiconque lit la
base --- ou vole la clé du coffre --- obtient l'accès à toutes les machines.

Avec une clé, le serveur ne détient que la \textbf{moitié publique}, qui ne
sert à rien à un attaquant. Compromettre la base de données d'InsaCloud ne
donne alors \emph{aucun} accès aux machines. Le compromis d'un système ne se
propage plus à l'autre.
\end{notion}

\limite
Trois limites, toutes à connaître :
\begin{itemize}[leftmargin=14pt,itemsep=2pt]
\item elle ne vaut que pour les machines louées \textbf{ensuite} --- la clé est
      injectée à la création du conteneur ;
\item le mot de passe \cmd{root} existe toujours sur la machine, car le
      terminal web et le bureau en ont besoin : la clé ferme la porte SSH, pas
      toutes les portes ;
\item une clé privée sans phrase de passe vaut ce que vaut la session de
      l'utilisateur sur son poste.
\end{itemize}

% -----------------------------------------------------------------------------
\mesure{Le terminal web passe par \cmd{login}}{mesure 16 · \fichier{entrypoint.sh}}

\menace
Un terminal dans le navigateur, c'est un shell exposé en HTTP. Mal protégé,
c'est un shell \cmd{root} offert à qui connaît l'URL.

\parade
\cmd{ttyd} ne lance pas un shell : il lance \cmd{/bin/login}, le programme de
connexion du système. L'utilisateur voit l'invite \cmd{login:} puis
\cmd{Password:}, exactement comme sur une console physique.

\danslecode
\begin{lstlisting}[language=bash]
# L'authentification est faite par `login` DANS le terminal (invite "login:"
# puis "Password:", comme une console), et non par une boite de dialogue HTTP.
\end{lstlisting}

Trois bénéfices, et le dernier n'est pas évident :

\begin{enumerate}[leftmargin=16pt,itemsep=2pt]
\item l'authentification passe par \textbf{PAM} : temporisation entre échecs,
      expiration de l'invite, journalisation --- gratuitement ;
\item l'expérience est familière, donc l'utilisateur n'est pas tenté de
      contourner ;
\item \textbf{couper l'authentification SSH par mot de passe n'enferme personne
      dehors} : le terminal web continue de fonctionner, puisqu'il ne passe pas
      par SSH. L'utilisateur qui perd sa clé garde une porte.
\end{enumerate}

\begin{piege}{La version précédente utilisait Basic Auth, et c'était pire}
Protégé par une authentification HTTP basique, le terminal affichait une page
blanche dans plusieurs navigateurs --- la négociation se faisait avant
l'ouverture de la WebSocket. Le remplacement par \cmd{login} a corrigé le
défaut d'usage \emph{et} amélioré la sécurité : Basic Auth transmet le secret à
chaque requête, \cmd{login} une seule fois dans un canal déjà établi.
\end{piege}

\limite
\cmd{ttyd} sert le terminal en HTTP si aucun certificat n'est monté. Sur le
parc déployé, le rôle \cmd{tls\_cert} fournit un certificat par nœud et
\cmd{ttyd} passe en \cmd{-{}-ssl} ; en pile locale de démonstration, le trafic
est en clair sur \cmd{localhost}.

% -----------------------------------------------------------------------------
\mesure{Le bureau graphique}{mesure 17 · \fichier{entrypoint.sh}}

\menace
Un serveur VNC exposé est une cible classique : le protocole est ancien et son
authentification est faible.

\parade
Trois précautions : le port VNC n'est \textbf{jamais publié}, le mot de passe
du bureau est \textbf{distinct} de celui de \cmd{root}, et le pont vers le
navigateur passe en TLS dès qu'un certificat est disponible.

\danslecode
\begin{lstlisting}[language=bash]
Xvnc :1 -geometry 1280x800 -depth 24 -rfbport 5901 -localhost \
        -SecurityTypes VncAuth -PasswordFile /root/.vnc/passwd -AlwaysShared
websockify $NOVNC_TLS --web $NOVNC_DIR 6080 localhost:5901
\end{lstlisting}

\cmd{-localhost} est le point essentiel : le port 5901 n'est joignable que
depuis l'intérieur du conteneur. Le seul port publié est le 6080, celui de
noVNC. Le mot de passe du bureau est dérivé du mot de passe \cmd{root} mais
tronqué à huit caractères, et stocké dans un fichier en \cmd{0600}.

\begin{piege}{Huit caractères, et ce n'est pas un choix}
Le protocole VncAuth \textbf{tronque le mot de passe à huit caractères} : c'est
une limite du protocole, pas un oubli. Il a donc fallu exposer un mot de passe
de bureau distinct de celui de \cmd{root}, pour ne pas laisser croire à
l'utilisateur que son mot de passe de seize caractères protège aussi son
bureau.
\end{piege}

\limite
\textbf{VncAuth est faible} : un défi-réponse DES des années 1990. C'est
acceptable ici pour trois raisons cumulées --- la machine est éphémère, le mot
de passe n'est visible qu'en ouvrant le coffre pendant cinq minutes, et le
transport est chiffré dès qu'un certificat est monté. Sur une plateforme
permanente, il faudrait un proxy authentifié devant noVNC.
